\documentclass[12pt,titlepage]{report}

%% FORMATTING IMPORTS
\usepackage[margin=1.0in]{geometry}
\usepackage{hyperref}
\usepackage{graphicx}
\usepackage{amsmath}
\usepackage{amssymb}
\usepackage{algorithm}
\usepackage{algpseudocode}

\begin{document}
	\begin{titlepage}
		\begin{center}
			\vspace*{5cm}
			{\Huge \textbf{The RSOPT Project}} \\
			\vspace*{2cm}
			{\large
				\textit{prepared by} \\
				\vspace*{0.2cm}
				Joseph R. Hobbs \\
				\vspace*{0.3cm}
				\href{mailto:rsopt@linotype.pub}{\texttt{rsopt@linotype.pub}}
			} \\
			\vspace*{2cm}
			{\large
				\textit{last updated} \\
				\vspace*{0.2cm}
				\today
			}
			\vfill
		\end{center}
	\end{titlepage}
	
	\tableofcontents
	
	\part{Introduction}
	
	\chapter{About RSOPT}
	
	TODO
	
	\chapter{Convex Optimization}
	
	TODO

	\chapter{Software Concept}
	
	TODO
	
	\part{Using RSOPT}
	
	\chapter{Solver Configuration}
	
	TODO
	
	\part{Convex Programs}
	
	\chapter{Linear Programs}
	
	TODO
	
	\chapter{Quadratic Programs}
	
	This chapter covers various forms of \textbf{quadratic programs} supported by the RSOPT project.  This document defines quadratic programs as convex programs whose objective has a constant Hessian with no negative eigenvalues.
	
	\section{Unconstrained Quadratic Programs}
	
	\textbf{Unconstrained quadratic programs} (\texttt{UnconstrainedQuadraticProgram<N>}) are convex programs of the form
	
	\begin{equation}
		\min_{x \in \mathbb{R}^n} \frac12 x^\top A x + b^\top x
	\end{equation}
	
	where \( A \in \mathbb{R}^{n \times n} \), \( A \succeq 0 \), and \( b \in \mathbb{R}^n \).  The relevant stationarity condition is
	
	\begin{equation}
		Ax^\star + b = 0 \Leftrightarrow x^\star = -A^{-1} b.
	\end{equation}
	
	We solve unconstrained quadratic programs using the \textbf{conjugate gradient} (CG) method.
	
	\paragraph{Definition.} Two vectors \( p_i \) and \( p_j \) are \textit{conjugate} with respect to \( A \) if \( p_i^\top A p_j = 0 \).
	\vspace{0.4cm}
	
	CG solves an \( n \)-dimensional system in at most \( n \) iterations by selecting conjugate search directions.  The gradient of the objective at iterate \( x_k \) is
	
	\begin{equation}
		\nabla J_k = A x_k + b .
	\end{equation}
	
	CG selects \( x_0 := 0 \) and \( p_0 := \nabla J_0 = b \) as the first search direction.  Subsequent search directions \( p_k \) (\( k > 0 \)) must be conjugate with all preceding search directions with respect to \( A \).  We write the recurrence as
	
	\begin{equation}
		p_{k+1} = A x_{k+1} + b + \sum_{i=0}^k \beta_i p_i
	\end{equation}
	
	and constrain \( \beta_i \) using the definition of conjugate vectors.  For all \( j = 0, 1, \cdots k \), we have
	
	\begin{equation}
		(A x_{k+1} + b + \sum_{i=0}^k \beta_i p_i)^\top A p_j = 0.
	\end{equation}
	
	Expanding gives us
	
	\begin{equation}
		x_{k+1}^\top A^\top A p_j + b^\top A p_j + \sum_{i=0}^k \beta_i p_i^\top A p_j = 0 .
	\end{equation}
	
	Notice that, by our ``inductive hypothesis", \( p_0, p_1, \cdots p_{k-1} \) are already pairwise conjugate with respect to \( A \).  We can now focus only on the case \( j = k \) and take all other \( \beta_i = 0 \).
	
	\begin{equation}
		x_{k+1}^\top A^\top A p_k + b^\top A p_k + \beta_k p_k^\top A p_k = 0 .
	\end{equation}
	
	This gives us a closed form for \( \beta_k \).
	
	\begin{equation}
		\beta_k = -\frac{(\nabla J_{k+1})^\top A p_k}{p_k^\top A p_k}
	\end{equation}
	
	Therefore, the step direction \( p_k \) for \( k > 0 \) is given by
	
	\begin{equation}
		\boxed{p_{k+1} = \nabla J_{k+1} - \frac{(\nabla J_{k+1})^\top A p_k}{p_k^\top A p_k} p_k. }
	\end{equation}
	
	After selecting the step direction, it is necessary to select the step size \( \alpha_k \).  We select \( \alpha_k \) according to the QP
	
	\begin{equation}
		\min_{\alpha_k \in \mathbb{R}} \frac12 (x_k + \alpha_k p_k)^\top A (x_k + \alpha_k p_k) + b^\top (x_k + \alpha_k p_k)
	\end{equation}
	
	Expanding the objective and dropping all terms not dependent on \( \alpha \), we have
	
	\begin{equation}
		\frac{d}{d\alpha_k} \left[ \alpha_k p_k^\top A x_k + \frac12 \alpha_k^2 p_k^\top A p_k + \alpha_k b^\top p_k \right] = 0.
	\end{equation}
	
	We find that \( \alpha_k \) is also determined analytically as
	
	\begin{equation}
		\boxed{\alpha_k = \frac{-(\nabla J_k)^\top p_k}{p_k^\top A p_k}. }
	\end{equation}
	
	The complete CG algorithm is detailed in Algorithm \ref{alg:cg}.  For two conjugate directions \( p_i \) and \( p_j \) with respect to \( A \), we can interpret this algorithm in terms of the eigendecomposition of \( A \).
	
	\begin{equation}
		p_i^\top A p_j = p_i^\top S \Lambda S^\top p_j
	\end{equation}
	
	We can interpret the algorithm as akin to taking orthogonal steps to minimize an objective of the form \( x^\top \Lambda x \), where \( \Lambda \) is a diagonal matrix containing the eigenvalues of \( A \).
	
	\begin{algorithm}[htbp!]
		\caption{Conjugate Gradient}
		\label{alg:cg}
		\begin{algorithmic}[1]
			\Require \( A, b \)
			\State \( x_0 \gets 0 \)
			\State \( p_0 \gets b \)
			\State \( g_0 \gets b \)
			\State \( \alpha_0 \gets -g_0^\top p_0 / p_0^\top A p_0 \)
			\State \( x_1 \gets x_0 + \alpha_0 p_0 \)
			\For{\( k = 1 \cdots n-1 \)}
			\State \( g_k \gets Ax_k + b \)
			\State \( p_k \gets g_k - (g_k^\top A p_{k-1}/p_{k-1}^\top A p_{k-1}) p_{k-1} \)
			\State \( \alpha_k \gets -g_k^\top p_k / p_k^\top A p_k \)
			\State \( x_{k+1} \gets x_k + \alpha_k p_k \)
			\EndFor \\
			\Return \( x_n \)
		\end{algorithmic}
	\end{algorithm}
	
	\section{Equality-Constrained Quadratic Programs}
	
	TODO
	
	\section{Inequality-Constrained Quadratic Programs}
	
	TODO
\end{document}